% Separate analytical increment. No supplied file is modified.
% Unique prefix: inc:AN03:ret:v1:
\documentclass[11pt]{article}
\usepackage[T1]{fontenc}
\usepackage[utf8]{inputenc}
\usepackage{lmodern}
\usepackage[margin=0.9in]{geometry}
\usepackage{amsmath,amssymb,amsthm,mathtools,booktabs,array,microtype}
\usepackage{xurl}
\usepackage[hidelinks]{hyperref}
\usepackage{fancyhdr}
\hypersetup{pdftitle={AN03: Two-sided nonlinear passage and Q2 ancestry retention},pdfauthor={Research proof increment for independent review}}
\pagestyle{fancy}\fancyhf{}
\fancyhead[L]{\small AN02/AN03: retention and two-sided passage}
\fancyhead[R]{\small Review increment v1}
\fancyfoot[C]{\thepage}
\setlength{\headheight}{14pt}
\setlength{\emergencystretch}{3em}
\allowdisplaybreaks[2]
\numberwithin{equation}{section}
\newtheorem{theorem}{Theorem}[section]
\newtheorem{lemma}[theorem]{Lemma}
\newtheorem{proposition}[theorem]{Proposition}
\newtheorem{corollary}[theorem]{Corollary}
\theoremstyle{remark}\newtheorem{remark}[theorem]{Remark}
\newcommand{\E}{\mathbb E}
\newcommand{\R}{\mathbb R}
\newcommand{\dd}{\,d}
\newcommand{\1}{\mathbf 1}
\newcommand{\scope}[1]{\par\noindent\textbf{Scope.} #1\par\smallskip}
\newcommand{\gap}[1]{\par\noindent\textbf{Open obligation.} #1\par\smallskip}
\title{Two-sided nonlinear passage and Q2 ancestry retention\\
\large A mean-feedback budget, corrected entry clocks, and an arrival--retention distinction}
\author{Research proof increment for independent review}
\date{Version 1 --- October 3, 2026}
\begin{document}
\maketitle
\begin{abstract}
For the original finite-noise initialized flow, the reviewed weighted-ancestry
bound gives Q2 ancestral mass at least a constant times
$S_0^{1-\lambda}$ at the early envelope time, including its exact
$q^2\log(1/S_0)$ correction. We express its conversion to captured mass
through a precise weighted net-retention functional. We also prove an
upper bound for Q2 labels arriving in the AN02 fixed sector: each such
label has only $O(S_0)$ mass when $q^2T$ is bounded. Thus growing Q2
ancestry and a late-arriving Q2 cohort are different quantities.
For the full leading population we derive upper- and lower-side pair
identities, retaining both the receiver term and the full donor field.
The lower-side cancellation yields an all-label $\ell^1$ distance bound,
without co-ordering. An exact nonlinear comparison closes this bound
through weak half-mass in terms of the time integral of the positive
barycentric imbalance. If the entry shape and retained weak seed have
powers $S_0^{1/2-\alpha}$ and $S_0^{1-\lambda}$, respectively, the resulting
conditional shape exponent is $\lambda/2-\alpha>0$ for every fixed
$0<\rho<1$. No initialized retention law, mean-feedback budget, full
resident passage to the selected orbit, or finite-noise transfer is
assumed proved. The reported sweep is motivation, not a proof premise.
\end{abstract}

\section{Contract and notation}\label{inc:AN03:ret:v1:contract}
The original-flow inputs are foundations P4--P5, the checkpoint labels
\texttt{pa:earlybounds} and \texttt{pa:targetcomparison}, and the reviewed
Theorem 2.1 of \path{AN02_strong_ancestry_and_tail_transport_v1.tex}.
Sections 3--5 use the \emph{defined leading system}
\texttt{pa:leading}/\texttt{pa:masslaws} and the exact overlap kernel
\texttt{pa:overlap}. Its relevant equations are restated. The comparison
characteristic below evolves in the \emph{same population field}, not
in the field of a fitted or coherent reference. The selected-orbit witness
target is the reviewed \path{AN06_learned_orbit_witnesses_v2.tex}.
All original sources and reviewed increments remain unchanged.

Write $\lambda=\rho^2\in(0,1)$, normalized time $\tau$, and initial-label
law $\lambda_0(\dd\alpha_0)=\dd\alpha_0/(2\pi)$. The symbol
$\alpha_{\rm sh}=\lambda\Phi(\rho a_\rho)/2$ denotes the checkpoint's
shape rate $\alpha$, not an initial label. As usual,
\[
 K(z)=\phi(z)+z\Phi(z),\qquad
 F(t,r)=\phi(\min(t,r))+r\Phi(\min(t,r)),\qquad
 a_\rho=\rho K(\rho a_\rho).
\]
The unique finite $a_\rho>0$ exists for every fixed $0<\rho<1$.
All exponents in this note refer to the initialization \emph{mass} $S_0$.

For the original flow, $P=(P_1+P_2)/2$ with
$P_1=\mathcal N(e_1,q^2I)$ and $P_2=\mathcal N(\rho e_2,q^2I)$.
The label vectors start at
$U_0=W_0=\sqrt{S_0}a_{\alpha_0}$. With $m=|U|^2=|W|^2$,
\[
 (\log m)'=r_{\rm rad}:=2b^\top R_f(\theta)a,
 \quad R_f(\theta)=\E_P[(X-f(X))X^\top\1_{\{a^\top X>0\}}].
\]
Put $a_*=1/2+q^2$ and
\[
 \overline S(\tau)=S_0e^{(1+2q^2)\tau},\quad
 \epsilon_{\rm al}=\overline S-S_0,\quad
 H(\tau)=S_0(2e^{2a_*\tau}-3e^{a_*\tau}+1).
\]
The exact target-only flow $\widehat U$ starts at the same vectors and has
$\widehat U'=\E_P[X(\widehat U^\top X)_+]$.
On $\epsilon_{\rm al}<\pi$, the reviewed comparison gives
\begin{equation}\label{inc:AN03:ret:v1:earlycomparison}
 |\theta-\widehat\theta|\le H,
 \qquad |\log(m/\widehat m)|\le2H,
 \qquad
 m_\tau(\alpha_0)\ge e^{-2H}S_0
 [e^\tau(\cos\alpha_0)_+^2+e^{\lambda\tau}(\sin\alpha_0)_+^2].
\end{equation}
This comparison is not continued through order-one strong learning here.

\section{Growing Q2 ancestry is not a late-arriving Q2 cohort}
\label{inc:AN03:ret:v1:ancestry}
\scope{Original finite-noise flow on its proved early interval, followed
by exact conditional retention inequalities. No captured-family law is
inferred from the name of an ancestral quadrant.}
Let $Q_2=(\pi/2,\pi)$ in initial-label space, and choose
$0<S_0<\delta\le1/100$. Define
\[
 T_\delta=\frac{\log(\delta/S_0)}{1+2q^2}.
\]
Then $S\le\delta$, $H(T_\delta)\le2\delta$, and
$\epsilon_{\rm al}<\pi$ throughout this interval.

\begin{proposition}[Q2 ancestral amplification at the early clock]
\label{inc:AN03:ret:v1:Q2mass}
For every fixed measurable $C\subset Q_2$,
\begin{equation}\label{inc:AN03:ret:v1:Q2weighted}
 \int_Cm_{T_\delta}\dd\lambda_0
 \ge e^{-4\delta}\delta^\lambda S_0^{1-\lambda}
 e^{-2\lambda q^2T_\delta}
 \int_C\sin^2\alpha_0\dd\lambda_0.
\end{equation}
In particular the total mass carried by Q2 ancestry satisfies
\begin{equation}\label{inc:AN03:ret:v1:Q2total}
 \nu^{\rm lab}_{T_\delta}(Q_2)
 \ge \frac{e^{-4\delta}\delta^\lambda}{8}
 S_0^{1-\lambda}e^{-2\lambda q^2T_\delta}.
\end{equation}
\end{proposition}
\begin{proof}
Use the positive second-coordinate term in
\eqref{inc:AN03:ret:v1:earlycomparison} and the exact identity
$S_0e^{\lambda T_\delta}
=\delta^\lambda S_0^{1-\lambda}e^{-2\lambda q^2T_\delta}$.
The integral of $\sin^2\alpha_0$ over Q2 in $\lambda_0$ is $1/8$.
\end{proof}
This is a lower bound, not a matching asymptotic equivalence. The labels
may have crossed toward the strong axis. Their current weak-chart mass
and their useful weak response are not identified by
\eqref{inc:AN03:ret:v1:Q2total}.

\subsection{An upper bound on the AN02 arriving cohort}
Let $\varphi_T$ be the exact target-only angular flow and set
\[
 I=[2\pi/3,3\pi/4],\qquad C_T=\varphi_{-T}(I),\qquad
 C_T^{(2)}=C_T\cap Q_2.
\]
These are precisely analytical backward-sector cohorts; no trained
trajectory is used to choose them.

\begin{theorem}[No uniform weak-growth gain for fixed-sector arrivals]
\label{inc:AN03:ret:v1:arrivalupper}
For $T\ge0$ and every initial label in $C_T$,
\begin{equation}\label{inc:AN03:ret:v1:arrivalmasshat}
 \widehat m_T(\alpha_0)
 \le4S_0\cos^2\alpha_0\,e^{3q^2T/2}
 \le4S_0e^{3q^2T/2}.
\end{equation}
For $\alpha_0\in C_T^{(2)}$,
\begin{equation}\label{inc:AN03:ret:v1:arrivalangle}
 \tan(\pi-\alpha_0)
 \le\sqrt3\,e^{-(\lambda/2-3q^2/4)T}.
\end{equation}
Consequently,
\begin{equation}\label{inc:AN03:ret:v1:arrivalmeasure}
 \lambda_0(C_T^{(2)})
 \le\frac{\sqrt3}{2\pi}e^{-(\lambda/2-3q^2/4)T}.
\end{equation}
If $\epsilon_{\rm al}(T)<\pi$, the original-flow mass obeys
\begin{align}
 m_T(\alpha_0)&\le4e^{2H(T)}S_0e^{3q^2T/2}
 &&(\alpha_0\in C_T),\label{inc:AN03:ret:v1:arrivalmass}\\
 \int_{C_T^{(2)}}m_T\dd\lambda_0
 &\le\frac{2\sqrt3}{\pi}e^{2H(T)}S_0
 e^{-\lambda T/2+9q^2T/4}.\label{inc:AN03:ret:v1:arrivalQ2mass}
\end{align}
\end{theorem}
\begin{proof}
The exact coordinate formula, with $r=|\widehat U|$ and
$\widehat a=\widehat U/r$, is
\begin{equation}\label{inc:AN03:ret:v1:coordinatefield}
 \widehat U_j'=\frac{\mu_jqr}{2}K(\mu_j\widehat a_j/q)
 +d_q(\widehat a)\widehat U_j,
 \quad
 d_q=\frac{q^2}{2}[\Phi(\widehat a_1/q)+\Phi(\rho\widehat a_2/q)],
 \quad (\mu_1,\mu_2)=(1,\rho).
\end{equation}
The positive-coordinate half-spaces are forward invariant: at a zero
coordinate its derivative is strictly positive. The radius stays
positive because $(r^2)'=2\E_P(\widehat U^\top X)_+^2\ge0$.
For a label ending in $I$, $\widehat U_1(T)<0$, hence
$\widehat U_1(t)<0$ for every $0\le t\le T$. On this path
$d_q\le3q^2/4$. For $w=-\widehat U_1>0$,
\[
 w'=-\frac{qr}{2}K(\widehat a_1/q)+d_qw
 \le\frac{3q^2}{4}w,
 \qquad w(T)\le\sqrt{S_0}|\cos\alpha_0|e^{3q^2T/4}.
\]
Since $|\cos\widehat\theta(T)|\ge1/2$, one has $r(T)\le2w(T)$,
which proves \eqref{inc:AN03:ret:v1:arrivalmasshat}.
For a Q2 label, the positive second coordinate satisfies
$\widehat U_2(T)\ge\sqrt{S_0}\sin\alpha_0 e^{\lambda T/2}$.
At the endpoint $\widehat U_2(T)/w(T)\le\sqrt3$.
Combining these bounds proves \eqref{inc:AN03:ret:v1:arrivalangle}.
Thus $C_T^{(2)}$ is contained in the terminal initial-label interval
of length at most
$\arctan(\sqrt3 e^{-(\lambda/2-3q^2/4)T})$, which is at most its
argument. Divide by $2\pi$ to get \eqref{inc:AN03:ret:v1:arrivalmeasure}.
The logarithmic comparison transfers the pointwise upper bound; multiplying
by the label-measure bound proves \eqref{inc:AN03:ret:v1:arrivalQ2mass}.
No initial-to-final angle inverse for the original flow is used.
\end{proof}

For fixed $\delta$ and $\rho$, at $T=T_\delta$ under
$q^2T_\delta\to0$, the last bound is
\begin{equation}\label{inc:AN03:ret:v1:arrivalscale}
 \int_{C_{T_\delta}^{(2)}}m_{T_\delta}\dd\lambda_0
 =O_\delta(S_0^{1+\lambda/2}).
\end{equation}
Indeed,
$e^{-\lambda T_\delta/2+9q^2T_\delta/4}
=(S_0/\delta)^{\lambda/2}
 e^{(\lambda+9/4)q^2T_\delta}$.
The equality symbol in \eqref{inc:AN03:ret:v1:arrivalscale} denotes an
upper-order statement, not an asymptotic equivalence. Compared with the
lower bound \eqref{inc:AN03:ret:v1:Q2total}, this arriving subset carries
at most $O_\delta(S_0^{3\lambda/2})$ of total Q2 ancestral mass.
The reviewed ancestry formula is therefore valuable for \emph{retention
of earlier Q2 labels}, not a uniform $e^{\lambda T}$ improvement for
AN02's late-arriving sector cohort. Its initial $\sin^2\alpha_0$ weight
shrinks with $T$ for that cohort.

\subsection{A precise original-flow retention target}
Fix a later deterministic time $\tau_e\ge T_\delta$ and a prescribed
weak-chart set $\mathcal W_R$, defined by
$\theta=\pi/2-qu$, $\psi=\pi/2+qv$, $|u|,|v|\le R$ (modulo $2\pi$).
Define, using the true characteristic path,
\begin{align}
 A_{\alpha_0}&=\int_{T_\delta}^{\tau_e}r_{\rm rad}(t,\alpha_0)\dd t,
 \qquad
 \chi_{\alpha_0}=\min\{1,e^{A_{\alpha_0}}\},\nonumber\\
 \mathcal R_2(\tau_e)
 &=\int_{Q_2}\1_{\{(\theta,\psi)(\tau_e,\alpha_0)\in\mathcal W_R\}}
 \sin^2\alpha_0\,\chi_{\alpha_0}\dd\lambda_0.
 \label{inc:AN03:ret:v1:retentionfunctional}
\end{align}
Here $0\le\mathcal R_2\le1/8$. The net multiplier permits mass loss
followed by recovery; it does not charge every temporary loss separately.

\begin{proposition}[Exact ancestry-to-capture inequality]
\label{inc:AN03:ret:v1:retention}
Let $N_2^{\rm cap}(\tau_e)$ be the true mass in $\mathcal W_R$ with Q2
ancestry. Then
\begin{equation}\label{inc:AN03:ret:v1:retentionbound}
 N_2^{\rm cap}(\tau_e)
 \ge e^{-4\delta}\delta^\lambda S_0^{1-\lambda}
 e^{-2\lambda q^2T_\delta}\,\mathcal R_2(\tau_e).
\end{equation}
In particular, if a measurably specified initial-label set $C\subset Q_2$
is inside $\mathcal W_R$ at $\tau_e$ and satisfies
\[
 \int_C\sin^2\alpha_0\dd\lambda_0\ge w_2,
 \qquad A_{\alpha_0}\ge-B_2\quad(\alpha_0\in C),\qquad B_2\ge0,
\]
then $\mathcal R_2\ge w_2e^{-B_2}$.
\end{proposition}
\begin{proof}
The exact radial equation gives
$m_{\tau_e}=m_{T_\delta}e^{A_{\alpha_0}}
\ge m_{T_\delta}\chi_{\alpha_0}$.
Apply the pointwise bound underlying
\eqref{inc:AN03:ret:v1:Q2weighted} and integrate with the displayed
fixed-time chart indicator. No derivative of that indicator is taken,
so no moving-mask flux is omitted. The sufficient condition is immediate.
\end{proof}

There is a crude unconditional bound on the radial part. Loss dissipation
and Cauchy--Schwarz give
\begin{equation}\label{inc:AN03:ret:v1:radialfloor}
 \|R_f(\theta)\|\le C_E,
 \quad C_E=\left|1-\frac{S_0}{4}\right|
 \sqrt{\left(\frac12+q^2\right)
       \left(\frac{1+\lambda}{2}+2q^2\right)},
 \quad A_{\alpha_0}\ge-2C_E(\tau_e-T_\delta).
\end{equation}
For the matrix bound, pair $R_f$ with two unit vectors, use
$\E\|X-f(X)\|^2\le2\mathcal L(0)$ and
$\E(z^\top X)^2\le a_*$; here
$2\mathcal L(0)=(1-S_0/4)^2\operatorname{tr}\E XX^\top$.
Thus a delay $O(\log(1/q))$ has at worst a polynomial radial-retention
cost on bounded parameter sets. A delay proportional to $\log(1/S_0)$
can change the initialization exponent under this crude bound.

\gap{Neither a lower bound on $\mathcal R_2$, nor the chart membership
of a set with nonvanishing weighted ancestry, is proved here. The capture
problem is to prove such a bound through transit and return, while
controlling the competing population. Even if the mass is captured once,
subsequent chart persistence and recruitment flux remain separate tasks.}

\section{Exact two-sided pair identities in the full leading field}
\label{inc:AN03:ret:v1:two-sided}
\scope{The full, self-consistent leading population on a finite
chart-valid interval. No population ordering is assumed in the final
all-label bound.}
Let $0<M\le1$, $0\le N\le1$, $\varepsilon_M=1-M$, and
\[
 M'=M(1-M),\qquad N'=\lambda N(1-N),\qquad
 Y=\int y\dd\mu_s,\quad V=\int v\dd\mu_w,\quad K_w=\int K(u)\dd\mu_w.
\]
The exact leading strong characteristic field is
\begin{align}
 2f_x(x,y)&=-\varepsilon_Mx+\rho\phi(\rho x)-\rho\mathcal S(x)
 +\lambda[V+(1-N)y]\Phi(\rho x),\label{inc:AN03:ret:v1:field}\\
 2f_y(x,y)&=-\varepsilon_My-Y-K_w+\rho(1-N)K(\rho x),\nonumber\\
 \mathcal S(x)&=\int F(\rho x,\rho\widetilde x)\dd\mu_s.
 \nonumber
\end{align}
Use bounded-support characteristic solutions, with fixed family labels
and their normalized weights. Let $(c,d)$ be any characteristic of this
\emph{same} time-dependent field. For example, initialize it at the strong
barycentre; it need not carry positive mass. Define
\begin{align}
 X_s&=M^{-1}\int x\dd\mu_s,\qquad Y_s=M^{-1}\int y\dd\mu_s,\nonumber\\
 \Pi(r)&=\int(\widetilde x-r)_+\dd\mu_s,\qquad
 \Lambda_-(r)=\int(r-\widetilde x)_+\dd\mu_s,\nonumber\\
 U_c&=V+(1-N)d-c,\qquad U_c^-=U_c+M(c-X_s),\nonumber\\
 \mathcal B&=V+(1-N)Y_s-X_s,\qquad
 c_\rho=\frac{\rho^3\phi(0)}2.
 \label{inc:AN03:ret:v1:budgets}
\end{align}
The total weak moment $V$ is not a normalized weak mean. The notation
$\mathcal B$ here is the barycentric imbalance, not the passive-diagonal
average from AN02.

\begin{lemma}[Derivative and donor cancellation identities]
\label{inc:AN03:ret:v1:jacobian}
Writing $P(r)=\Phi(\rho r)$, one has exactly
\begin{align}
 \mathcal S'(r)&=\rho^2\phi(\rho r)\Pi(r),\nonumber\\
 2\partial_xf_x(r,d)&=-\varepsilon_M+
 \rho^3[V+(1-N)d-r-\Pi(r)]\phi(\rho r),\label{inc:AN03:ret:v1:dx}\\
 2\partial_yf_x(r,d)&=2\partial_xf_y(r,d)=\lambda(1-N)P(r),\nonumber\\
 \Pi(r)-\Lambda_-(r)&=M(X_s-r).\label{inc:AN03:ret:v1:donorcancel}
\end{align}
\end{lemma}
\begin{proof}
Differentiate the explicit $C^1$ overlap kernel:
$\partial_tF(t,z)=(z-t)_+\phi(t)$.
Its derivative is bounded on bounded sets, which justifies differentiating
under the finite measure. Use $\phi'(\rho r)=-\rho r\phi(\rho r)$
and $K'=\Phi$ to differentiate \eqref{inc:AN03:ret:v1:field}.
Finally $(z-r)_+-(r-z)_+=z-r$ proves
\eqref{inc:AN03:ret:v1:donorcancel} by integration.
\end{proof}

For $h>0$ define
\[
 \overline P_+(h)=\frac1h\int_c^{c+h}P(r)\dd r,
 \quad P_+=P(c+h),\qquad
 \overline P_-(h)=\frac1h\int_{c-h}^{c}P(r)\dd r,
 \quad P_-=P(c-h).
\]
At $h=0$ use their continuous values $P(c)$.

\begin{proposition}[Receiver-preserving upper and lower identities]
\label{inc:AN03:ret:v1:averagedpair}
For an upper co-ordered pair $x=c+h$, $y=d+k$, $h,k\ge0$, put $D=h+k$.
Then
\begin{align}
 D'={}&\left[-\frac{\varepsilon_M}{2}
       +\frac\lambda2(1-N)\overline P_+\right]D\nonumber\\
 &+\frac\lambda2[(1-N)k-h](P_+-\overline P_+)
 +\frac{\rho^3}{2}\int_c^{c+h}[U_c-\Pi(r)]\phi(\rho r)\dd r.
 \label{inc:AN03:ret:v1:upperexact}
\end{align}
For a lower co-ordered pair $x=c-h$, $y=d-k$, $h,k\ge0$, one has
\begin{align}
 D'={}&\left[-\frac{\varepsilon_M}{2}
       +\frac\lambda2(1-N)\overline P_-\right]D\nonumber\\
 &+\frac\lambda2[\varepsilon_M h-(1-N)k]
       (\overline P_--P_-)
 +\frac{\rho^3}{2}\int_{c-h}^{c}[U_c^--\Lambda_-(r)]\phi(\rho r)\dd r.
 \label{inc:AN03:ret:v1:lowerexact}
\end{align}
The identities apply on any time interval on which the indicated order
holds; leading same-field cooperativity preserves an initial such order.
\end{proposition}
\begin{proof}
For the upper pair, the fundamental theorem of calculus and
\eqref{inc:AN03:ret:v1:dx} give
\begin{align*}
 2h'&=-\varepsilon_Mh+
 \rho^3\int_c^{c+h}[U_c-(r-c)-\Pi(r)]\phi(\rho r)\dd r
 +\lambda(1-N)kP_+,\\
 2k'&=-\varepsilon_Mk+\lambda(1-N)h\overline P_+.
\end{align*}
Integration by parts gives
$\rho\int_c^{c+h}(r-c)\phi(\rho r)\dd r
=h(P_+-\overline P_+)$.
Substitute, add, and collect the coefficient of $D$ to obtain
\eqref{inc:AN03:ret:v1:upperexact}.
For the lower pair the corresponding first integral is
$\rho^3\int_{c-h}^{c}[U_c+(c-r)-\Pi(r)]\phi(\rho r)\dd r$,
with cross term $\lambda(1-N)kP_-$. The donor identity rewrites its
bracket as
\[
 U_c+(c-r)-\Pi(r)=U_c^-+\varepsilon_M(c-r)-\Lambda_-(r).
\]
Now
$\rho\int_{c-h}^{c}(c-r)\phi(\rho r)\dd r
=h(\overline P_--P_-)$.
Adding the two component equations proves
\eqref{inc:AN03:ret:v1:lowerexact}.
The vanishing-$h$ cases follow by continuity. The cross derivatives in
Lemma~\ref{inc:AN03:ret:v1:jacobian} are nonnegative for $N\le1$,
which gives same-field order by a first-contact argument.
\end{proof}

On the upper diagonal $k=h$ at $N=0$, the mismatch term in
\eqref{inc:AN03:ret:v1:upperexact} vanishes. In the resident field with
$c=d=a_\rho$, $\mu_s=M\delta_{(a_\rho,a_\rho)}$, $N=0$, its integral
term also vanishes for the upper pair. The rate becomes
$-\varepsilon_M/2+\lambda\overline P_+/2$, exactly the AN02 nonlinear
average. On the lower side, at $M=1$ the mismatch term is nonpositive
for every $k\ge0$. Dropping the donor term \emph{before} using
\eqref{inc:AN03:ret:v1:donorcancel} would lose this cancellation.

\section{An all-label distance bound and nonlinear half-mass passage}
\label{inc:AN03:ret:v1:passage}
For an arbitrary strong label, regardless of order, let
\[
 D_{\alpha_0}=|x_{\alpha_0}-c|+|y_{\alpha_0}-d|,
 \quad
 R_p=\left(\frac1M\int D^p\dd\mu_s\right)^{1/p}\ (1\le p<\infty),
 \quad R_\infty=\operatorname*{ess\,sup}D.
\]
The integrals mean pushforwards of the fixed label law; no one-to-one
angle map is needed.

\begin{theorem}[Two-sided, mixed-order differential bound]
\label{inc:AN03:ret:v1:alllabel}
Every strong label satisfies, in upper Dini form or almost everywhere,
\begin{equation}\label{inc:AN03:ret:v1:alllabelineq}
 D^+D\le
 \left[-\frac{(1-\lambda)(1-M)}2+
       \frac\lambda2(1-N)+b_c\right]D,
 \qquad
 b_c=c_\rho\max\{0,U_c,U_c^-\}.
\end{equation}
Moreover
\begin{equation}\label{inc:AN03:ret:v1:barybudget}
 b_c\le c_\rho(\mathcal B_++R_1).
\end{equation}
In particular, when $N(\tau_0)>0$, a direct integrated bound is
\begin{equation}\label{inc:AN03:ret:v1:Gdirect}
 D(t)\le
 \left(\frac{M(\tau_0)}{M(t)}\right)^{(1-\lambda)/2}
 \sqrt{\frac{N(t)}{N(\tau_0)}}
 e^{\int_{\tau_0}^t b_c}\,D(\tau_0).
\end{equation}
\end{theorem}
\begin{proof}
Let $h=|x-c|$, $k=|y-d|$. If $x>c$, multiply the exact $x-c$
equation by its sign and bound its cross term by
$\lambda(1-N)kP_+$. The derivative of $|y-d|$ is at most
$[-\varepsilon_Mk+\lambda(1-N)h\overline P_+]/2$.
These are the same upper inequalities as for the co-ordered upper pair,
even when $y-d<0$. The unintegrated upper formula in the proof of
Proposition~\ref{inc:AN03:ret:v1:averagedpair} has
$-(r-c)-\Pi(r)\le0$. Hence
\[
 D^+D\le[-\varepsilon_M/2+\lambda(1-N)/2+c_\rho(U_c)_+]D.
\]
If $x<c$, the same absolute-value argument gives upper inequalities of
the lower-pair form. Apply \eqref{inc:AN03:ret:v1:lowerexact}, discard
$-\Lambda_-$, and use
\[
 [\varepsilon_Mh-(1-N)k](\overline P_--P_-)
 \le\varepsilon_Mh\le\varepsilon_MD,
 \quad 0\le\overline P_-\le1.
\]
This proves \eqref{inc:AN03:ret:v1:alllabelineq} on this side and also
bounds the upper side, because $1-\lambda\le1$.
At $x=c$ the upper derivative of $|x-c|$ is
$\lambda(1-N)P(c)k/2$, and at $y=d$ use the analogous absolute
value of the cross velocity. Thus the same bound holds at coordinate
equalities. At $D=0$ it holds by same-field uniqueness.

By Jensen,
$|c-X_s|+|d-Y_s|\le R_1$.
The exact identities
\[
 U_c=\mathcal B+(1-N)(d-Y_s)-(c-X_s),\qquad
 U_c^-=\mathcal B+(1-N)(d-Y_s)-(1-M)(c-X_s)
\]
prove \eqref{inc:AN03:ret:v1:barybudget}.
Finally integrate \eqref{inc:AN03:ret:v1:alllabelineq} and use the two
logistic identities $\int(1-M)=\log(M(t)/M(\tau_0))$ and
$\lambda\int(1-N)=\log(N(t)/N(\tau_0))$.
\end{proof}

The all-label bound is coarser on the lower side during strong learning
than AN02's upper estimate, but it retains the same global weak-stage
rate $\lambda/2$. It cannot be used backwards as an initialized
strong-contraction theorem. It does not require a pointwise small tail,
only the stated leading-system existence and moment hypotheses.

\begin{theorem}[Nonlinear passage to weak half-mass with a mean budget]
\label{inc:AN03:ret:v1:halfpassage}
Assume the leading solution and the same-field reference are defined on
$[\tau_0,\tau_h]$, where $0<N_0=N(\tau_0)<1/2$ and $N(\tau_h)=1/2$.
Define
\begin{align}
 L(t)&=\sqrt{\frac{N(t)}{N_0}}
       \exp\left(c_\rho\int_{\tau_0}^t\mathcal B_+(r)\dd r\right),
 \qquad L(\tau_0)=1,\label{inc:AN03:ret:v1:L}\\
 J_L(t)&=\int_{\tau_0}^t L(r)\dd r,
 \qquad Q(t)=1-c_\rho R_1(\tau_0)J_L(t).
 \label{inc:AN03:ret:v1:denominator}
\end{align}
Whenever $Q(t)>0$ throughout the interval,
\begin{equation}\label{inc:AN03:ret:v1:nonlinearG}
 D(t,\alpha_0)\le\frac{L(t)}{Q(t)}D(\tau_0,\alpha_0),
 \qquad
 R_p(t)\le\frac{L(t)}{Q(t)}R_p(\tau_0)
 \quad(1\le p\le\infty).
\end{equation}
Also
\begin{equation}\label{inc:AN03:ret:v1:integralL}
 J_L(t)\le\frac4\lambda[L(t)-1].
\end{equation}
Consequently, if
\begin{equation}\label{inc:AN03:ret:v1:smallness}
 \frac{4c_\rho}{\lambda}R_1(\tau_0)L(\tau_h)\le\frac12,
\end{equation}
then $Q\ge1/2$ and
\begin{equation}\label{inc:AN03:ret:v1:halfradius}
 R_p(\tau_h)\le
 2\,\frac{e^{c_\rho\int_{\tau_0}^{\tau_h}\mathcal B_+}}
          {\sqrt{2N_0}}R_p(\tau_0).
\end{equation}
The strong radius about its own barycentre is at most $2R_p$.
\end{theorem}
\begin{proof}
Drop the nonpositive strong-mass term in
\eqref{inc:AN03:ret:v1:alllabelineq} and use
\eqref{inc:AN03:ret:v1:barybudget}. With
$a(t)=\lambda(1-N)/2+c_\rho\mathcal B_+$, every label satisfies
$D^+D\le[a(t)+c_\rho R_1]D$.
The normalized strong weights are constant, so integration gives
\[
 R_1'\le a(t)R_1+c_\rho R_1^2
\]
almost everywhere. The explicit scalar solution with the same initial
value is $R_1(\tau_0)L(t)/Q(t)$, as verified by differentiation.
Scalar comparison on the interval of positive denominator bounds $R_1$
by this expression. This also covers zero initial radius by uniqueness.
It follows that
\[
 \int_{\tau_0}^t c_\rho R_1(r)\dd r
 \le-\log Q(t).
\]
Integrating the labelwise inequality proves
\eqref{inc:AN03:ret:v1:nonlinearG}; the $L^p$ bounds follow from the
constant normalized weights (the essential-supremum case is direct).

Since $N\le1/2$, $a(t)\ge\lambda/4$ and $L'=aL\ge\lambda L/4$.
Integrate this inequality to prove \eqref{inc:AN03:ret:v1:integralL}.
Both $L$ and $J_L$ are nondecreasing. Condition
\eqref{inc:AN03:ret:v1:smallness} therefore keeps $Q\ge1/2$ on the
whole interval and closes the comparison without a presumed small-shape
bootstrap. Finally the barycentre-reference distance is at most $R_1$;
Minkowski and $R_1\le R_p$ give the last claim.
\end{proof}

\paragraph{A usable sufficient mean budget.}
If
$\mathcal B_+\le C_B N+r_B$, $r_B\ge0$, and
$\int_{\tau_0}^{\tau_h}r_B\le B_{\rm err}$, then
\begin{equation}\label{inc:AN03:ret:v1:meanbudget}
 \int_{\tau_0}^{\tau_h}\mathcal B_+
 \le\frac{C_B}{\lambda}\log[2(1-N_0)]+B_{\rm err}
 \le\frac{C_B\log2}{\lambda}+B_{\rm err}.
\end{equation}
This follows by integrating $N'=\lambda N(1-N)$.
On the AN06 selected coherent branch before half-mass,
$\mathcal B=n(v-y)-(x-y)\le(8/5)n$ on its reviewed parameter interval.
Thus a bounded budget is analytically consistent with that branch.
The population has \emph{not} been proved to satisfy the same bound:
its mean tracking and weak moments must be controlled separately.

\begin{corollary}[An all-label first-exit profile]
\label{inc:AN03:ret:v1:exitprofile}
Let $G(t)=L(t)/Q(t)$ on a positive-denominator interval. Write
$\eta_0$ for the normalized fixed strong-label law and
$\mathcal P_0(r)=\eta_0\{D(\tau_0)\ge r\}$.
Then, for any $H_*>0$ and $T\le\tau_h$,
\begin{equation}\label{inc:AN03:ret:v1:exitmass}
 \eta_0\left\{\sup_{\tau_0\le t\le T}D(t)\ge H_*\right\}
 \le\mathcal P_0(H_*/G(T)).
\end{equation}
Multiplication by $M(T)$ gives the mass of those fixed labels at time
$T$. If $\mathcal P_0(r)\le C(\varepsilon_0/r)^p$, the right side is
at most $C(\varepsilon_0G(T)/H_*)^p$.
\end{corollary}
\begin{proof}
$L$ increases and $Q$ decreases, so $G$ is nondecreasing.
A first hit of $H_*$ implies $D(\tau_0)\ge H_*/G(T)$ by
\eqref{inc:AN03:ret:v1:nonlinearG}; labels already outside at entry
satisfy the same implication. Integrate this inclusion in the fixed
normalized label law. No time-dependent mask is differentiated.
\end{proof}
The theorem still requires leading-chart validity until $T$. It does not
authorize continuing an escaped label in a bounded physical chart or
ignoring that label's later residual contribution.

\section{Corrected entry clocks and the conditional exponent}
\label{inc:AN03:ret:v1:clocks}
\begin{proposition}[Entry-clock conversion with a defect]
\label{inc:AN03:ret:v1:clockconversion}
For an absolutely continuous family mass $0<N<1$ with
$N'=\lambda N(1-N)+r_N$,
\begin{equation}\label{inc:AN03:ret:v1:logit}
 \log\frac{N(\tau_e)}{1-N(\tau_e)}
 =\log\frac{N(0)}{1-N(0)}+\lambda\tau_e+
 \int_0^{\tau_e}\frac{r_N}{N(1-N)}\dd\tau.
\end{equation}
Suppose, at fixed $q$, that
\[
 \log\frac{N(0)}{1-N(0)}=-\gamma_0\log(1/S_0)+O_q(1),\quad
 \tau_e=\chi\log(1/S_0)+O_q(1),\quad
 \int_0^{\tau_e}\frac{r_N}{N(1-N)}\dd\tau=O_q(1).
\]
If $\gamma_e=\gamma_0-\lambda\chi>0$, then
$N(\tau_e)\asymp_q S_0^{\gamma_e}$.
For the closed logistic model with $\gamma_0=\chi=1$, this gives
$\gamma_e=1-\lambda$, not $1$.
\end{proposition}
\begin{proof}
Differentiate the logit and integrate. Under the hypotheses its endpoint
is $-\gamma_e\log(1/S_0)+O_q(1)$; since $\gamma_e>0$, the endpoint
mass tends to zero and its odds and mass are comparable.
\end{proof}
This conversion cannot be applied to an uncharted or changing family
without a mass-law/flux budget. The large relative-defect integral near
zero belongs here in the \emph{entry-law derivation}. Once the actual
half-mass time is used for the witness clock, the local AN06 mass defect
is integrated only on the bounded clock-to-witness hull; these are two
different uses of the same identity.

\begin{corollary}[Positive conditional global-rate shape exponent]
\label{inc:AN03:ret:v1:exponent}
Under the hypotheses of Theorem~\ref{inc:AN03:ret:v1:halfpassage}, suppose
for some $p\ge1$ that the \emph{entry} estimates are
\begin{equation}\label{inc:AN03:ret:v1:entryhyp}
 N_0\ge c_w(q)S_0^{1-\lambda},\qquad
 R_p(\tau_0)\le C_s(q)S_0^{1/2-\alpha_{\rm sh}},\qquad
 \int_{\tau_0}^{\tau_h}\mathcal B_+\le B_0(q).
\end{equation}
Set
\[
 A(q)=\frac{C_s(q)e^{c_\rho B_0(q)}}{\sqrt{2c_w(q)}},
 \qquad
 \eta_{\rm glob}=\frac\lambda2-\alpha_{\rm sh}
 =\frac\lambda2\Phi(-\rho a_\rho)>0.
\]
If $A(q)S_0^{\eta_{\rm glob}}\le\lambda/(8c_\rho)$, then
\begin{equation}\label{inc:AN03:ret:v1:exponentbound}
 R_p(\tau_h)\le2A(q)S_0^{\eta_{\rm glob}}.
\end{equation}
On $\rho\in[13/20,7/10]$, the reviewed bound
$\Phi(\rho a_\rho)\le71/100$ gives the explicit uniform inequality
\begin{equation}\label{inc:AN03:ret:v1:uniformeta}
 \eta_{\rm glob}\ge\frac{4901}{80000}>0.
\end{equation}
\end{corollary}
\begin{proof}
$R_1(\tau_0)\le R_p(\tau_0)$, and
$L(\tau_h)\le e^{c_\rho B_0}/\sqrt{2c_wS_0^{1-\lambda}}$.
Their product is at most $A(q)S_0^{\lambda/2-\alpha_{\rm sh}}$.
The specified inequality implies
\eqref{inc:AN03:ret:v1:smallness}; apply
\eqref{inc:AN03:ret:v1:halfradius}.
Positivity follows from the positive Gaussian tail at every finite
$\rho a_\rho$. Finally
$(\lambda/2)(1-71/100)\ge(169/800)(29/100)$.
\end{proof}

For comparison, the \emph{frozen resident} multiplier with the same entry
power has formal exponent
\begin{equation}\label{inc:AN03:ret:v1:frozeneta}
 \eta_{\rm frozen}
 =\frac12-\alpha_{\rm sh}-(1-\lambda)\frac{\alpha_{\rm sh}}\lambda
 =\frac12\Phi(-\rho a_\rho)>0.
\end{equation}
Equation \eqref{inc:AN03:ret:v1:frozeneta} is an algebraic identity for
that approximation, not a replacement of the moving population field.
The proved conditional result is \eqref{inc:AN03:ret:v1:exponentbound}.
No uniform positive exponent is claimed as $\rho\to0$ or $\rho\to1$.
The old $c_0=1$ balance remains the algebra for a different hypothetical
entry law; its root is not a proved phase transition of the initialized
population.

In a joint limit, the relevant requirement is
$\log A(q)=o(\log(1/S_0))$, uniformly on the chosen compact ratio
interval. Polynomial $q$-costs, including $B_0(q)=O(\log(1/q))$, are
absorbed by $\log(1/S_0)/\log(1/q)\to\infty$.
An unspecified $c_w(q)$ need not have this property. For example, a
retention cost $c_w(q)=e^{-c/q^2}$ is not absorbed under
$q^2\log(1/S_0)\to0$. The positive $S_0$ exponent does not eliminate
the quantitative retention problem.

\section{Empirical interpretation and the unresolved initialized link}
\label{inc:AN03:ret:v1:status}
The reported pilot sweep gives a weak mass exponent near $0.576$ at
strong half-mass and a strong input-spread exponent near $0.36$ at an
earlier learned strong threshold. These marginal input statistics do not
yet bound the joint input/output $R_p$ used here; percentile widths alone
also do not bound a high moment or a uniform radius. Those reported values motivate
\eqref{inc:AN03:ret:v1:entryhyp}; they were not independently recomputed
for this increment. The reported masks, single-run status, and finite
range must be retained. The formulas above neither require these data
nor show that the same scaling persists in the deeper proposed round-trip
regime. No $\rho=0.85$ sweep was run.

\paragraph{Recruitment comparison.}
The reviewed zero-noise target-only outer tail has power
$S_0^{3(1-\lambda)/2}q^{-3}$ at its specified comparison clock.
Relative to a \emph{hypothetical retained} seed $c_w(q)S_0^{1-\lambda}$,
its ratio has scale
\[
 c_w(q)^{-1}S_0^{(1-\lambda)/2}q^{-3}.
\]
At $\lambda=4/9$, the initialization-power advantage is $5/18$.
The heuristic threshold $\log(1/S_0)/\log(1/q)>54/5$ ignores any
additional $q$-power in $c_w$; it is exact exponent accounting only when
that omitted prefactor is controlled accordingly. If
$c_w(q)\gtrsim q^k$, the corresponding sufficient strict power condition
is $(1-\lambda)\log(1/S_0)/2>(3+k)\log(1/q)$ with a margin tending
to infinity. This remains a comparison of proposed source laws, not a
proof that recruited mass is negligible in the original flow.

\paragraph{What is improved.}
The new passage estimate covers upper, lower, and mixed-order strong
labels in one full-field bound. Its nonlinear closure uses $R_1$, the
mean-imbalance budget, and the weak mass growth; no additional
$\log(1/N_0)$ factor is introduced. The local averaged receiver formulas
remain available for sharper constants. Thus the earlier statement that
a local-core/exit decomposition is necessary for the \emph{sign} of the
power balance is not retained under the new candidate entry exponent.
Tail and outside-chart control are nevertheless still necessary for the
full finite-noise/witness argument.

\paragraph{What is not improved by changing the exponent alone.}
The same-field reference can follow the wrong mean trajectory; its
closeness to the selected coherent branch is not a conclusion of a
spread estimate. Neither the weak angular distribution nor a complete
AN06 distance follows from $R_p$ alone. The entry shape law, a useful
lower bound on $\mathcal R_2$, the mean-feedback budget
\eqref{inc:AN03:ret:v1:meanbudget}, and persistence/flux estimates for the
charted families remain to be proved from isotropic initialization.
The bridge-sized canonical distribution is not shown to meet the
smallness condition. Finite-noise rates, normalized weight drift,
uncharted populations, and probe-gate strips remain AN7 obligations.

The next initialized target is therefore explicit: prove, at an
analytically chosen learned entry time, a lower bound on the
\emph{weighted net retention} $\mathcal R_2$ with usable $q$-dependence,
together with a strong mean/shape estimate and a mean-imbalance integral
budget. The hypotheses in \eqref{inc:AN03:ret:v1:entryhyp} are a proved
sufficient passage interface, not an initialized theorem. Theorems A
and B remain open.

\begin{thebibliography}{9}
\raggedright
\bibitem{F} \path{RELU_POPULATION_FOUNDATIONS_v2.tex}.
Frozen supplied foundations, Gaussian and loss identities.
\bibitem{P} \path{THEOREM_A_ANALYTICAL_PROGRESS.tex}, version 1.0.
\texttt{pa:earlybounds}, \texttt{pa:targetcomparison}, \texttt{pa:overlap},
\texttt{pa:leading}, \texttt{pa:masslaws}, \texttt{pa:resident},
\texttt{pa:multipliers}.
\bibitem{A2} \path{AN02_strong_ancestry_and_tail_transport_v1.tex}.
Reviewed weighted ancestry, exact target-only outer law, and upper-tail
transport; the lower-side and all-label extensions here are new.
\bibitem{Arr} \path{AN02_weak_side_arrival_v1.tex}.
Reviewed target-only backward-sector cohort and original-flow arrival.
\bibitem{A6} \path{AN06_learned_orbit_witnesses_v2.tex}.
Reviewed learned limiting witnesses, half-mass clock, and small-shape
population tolerance.
\bibitem{Rev} User's current review and reported pilot $S_0$-sweep
measurements. Empirical motivation only; no trajectory data are a proof
premise.
\end{thebibliography}
\end{document}
